% Problem record op_5a17b19b8adeb191. This TeX file is derived from
% database/problems_json/op_5a17b19b8adeb191.json by scripts/sync-tex.mjs; edit the JSON record
% and rerun the script rather than editing this file.
\section{Smallest qubit number for a non-semi-Clifford third-level gate}
\label{sec:op_5a17b19b8adeb191}

\paragraph{Problem.}
What is the smallest number of qubits on which the third level of the
Clifford hierarchy contains a gate that is not semi-Clifford?  For $n\geq1$
qubits, let $\mathcal{P}_n$ be the Pauli group with arbitrary global phases,
consisting of the operators $e^{i\theta}P_1\otimes\cdots\otimes P_n$ with
$\theta\in\mathbb{R}$ and $P_j\in\{I,X,Y,Z\}$, and define the Clifford
hierarchy recursively by
\begin{equation}
  \mathcal{C}_1(n):=\mathcal{P}_n,
  \qquad
  \mathcal{C}_{k+1}(n):=
  \bigl\{U\in U(2^n):UPU^\dagger\in\mathcal{C}_k(n)
  \ \text{for every}\ P\in\mathcal{P}_n\bigr\},
  \qquad k\geq1.
  \label{eq:5a17-hierarchy}
\end{equation}
Thus $\mathcal{C}_2(n)$ is the Clifford group.  A unitary is \emph{semi-Clifford}
when it can be written as $C_LDC_R$ with $C_L,C_R\in\mathcal{C}_2(n)$ and $D$
diagonal in the computational basis; let $\mathrm{SC}(n)$ denote the set of such
unitaries.  The quantity sought is
\begin{equation}
  n_{\min}:=\min\bigl\{n\geq1:
  \mathcal{C}_3(n)\setminus\mathrm{SC}(n)\neq\varnothing\bigr\},
  \label{eq:5a17-nmin}
\end{equation}
where $\mathcal{C}_3(n)$ is the third level in Eq.~\eqref{eq:5a17-hierarchy}.
The known results give $n_{\min}\in\{5,6,7\}$.  A complete answer determines
the value of $n_{\min}$ in Eq.~\eqref{eq:5a17-nmin}; equivalently, it decides
for $n=5$ and for $n=6$ whether every gate in $\mathcal{C}_3(n)$ is
semi-Clifford.

\subsection*{Status}

Unsolved

\subsection*{Source}

Zeng, Chen, and Chuang conjectured that every third-level gate is
semi-Clifford and proved it for at most three qubits
\sourcecite{ref:5a17-zcc08}{ZCC08}; Beigi and Shor reported the seven-qubit
counterexample of Gottesman and Mochon \sourcecite{ref:5a17-bs10}{BS10}.
Anderson and Connelly identified four qubits as the smallest case in which the
structure of the third level was not fully known, and settled it
\sourcecite{ref:5a17-ac25}{AC25}.  He, Robitaille, and Tan state explicitly, in
Section~2.1 (page~2:6) of the TQC 2026 version, that it is an open problem
whether a third-level gate on five or six qubits can fail to be semi-Clifford
\sourcecite{ref:5a17-hrt26}{HRT26}.  Eq.~\eqref{eq:5a17-nmin} restates that
explicit question as the determination of the least such number of qubits.

\subsection*{Progress}

\begin{itemize}
  \item Zeng, Chen, and Chuang proved that every gate of every level is
  semi-Clifford on one or two qubits (Theorem~1) and that every third-level gate
  on three qubits is semi-Clifford (Theorem~2, by exhaustive computation), while
  for every level $k\geq4$ some three-qubit gates are not semi-Clifford
  (Theorem~3):
  \begin{equation}
    \mathcal{C}_3(n)\subseteq\mathrm{SC}(n)\quad(1\leq n\leq3),
    \qquad
    \mathcal{C}_k(3)\not\subseteq\mathrm{SC}(3)\quad(k\geq4).
    \label{eq:5a17-zcc}
  \end{equation}
  The left part of Eq.~\eqref{eq:5a17-zcc} gives $n_{\min}\geq4$.  Together with
  the two-qubit theorem, the right part shows that the analogue of $n_{\min}$
  at every level $k\geq4$ equals three, so the third level is the only level at
  which this threshold is undetermined \sourcecite{ref:5a17-zcc08}{ZCC08}.

  \item Beigi and Shor reported, in Section~1.1, a seven-qubit construction of
  Gottesman and Mochon: a product of three controlled-swap gates and four doubly
  controlled $Z$ gates that lies in $\mathcal{C}_3(7)$ but is not semi-Clifford.
  Hence
  \begin{equation}
    \mathcal{C}_3(7)\setminus\mathrm{SC}(7)\neq\varnothing,
    \qquad
    n_{\min}\leq7.
    \label{eq:5a17-seven}
  \end{equation}
  Their Theorem~1.1 shows that every third-level gate is generalized
  semi-Clifford \sourcecite{ref:5a17-bs10}{BS10}; by Proposition~2 of
  \sourcecite{ref:5a17-zcc08}{ZCC08}, such a gate factors as $C_L\Pi DC_R$ with
  Clifford gates $C_L,C_R$, a computational-basis permutation $\Pi$, and a
  diagonal $D$.  For a gate in $\mathcal{C}_3(n)\setminus\mathrm{SC}(n)$, the
  permutation $\Pi$ in such a factorization cannot be a Clifford gate.

  \item Anderson and Connelly proved, in Section~4 of a 2025 preprint, that
  every four-qubit third-level gate is semi-Clifford, so that, with
  Eqs.~\eqref{eq:5a17-zcc} and~\eqref{eq:5a17-seven},
  \begin{equation}
    \mathcal{C}_3(4)\subseteq\mathrm{SC}(4),
    \qquad
    n_{\min}\in\{5,6,7\}.
    \label{eq:5a17-four}
  \end{equation}
  The proof starts from the generalized semi-Clifford form, uses the
  classification of four-bit permutations into affine equivalence classes (only
  one class, containing the Toffoli gate, lies in the third level but not the
  Clifford group), reduces the diagonal factor to finitely many equivalence
  classes, and checks every remaining class by computer (Algorithm~1).  The
  authors remark that the approach might settle five qubits with ample
  computational resources, whereas six qubits likely requires further
  theoretical reductions \sourcecite{ref:5a17-ac25}{AC25}.

  \item He, Robitaille, and Tan characterized the permutation gates in the third
  level.  In Appendix~A of the full version they prove that every third-level
  permutation gate on at most six qubits is semi-Clifford, and that seven is the
  least number of qubits carrying a non-semi-Clifford third-level permutation
  gate (Theorem~A.6):
  \begin{equation}
    \mathcal{C}_3(n)\cap\mathrm{Perm}(2^n)\subseteq\mathrm{SC}(n)
    \qquad(n\leq6),
    \label{eq:5a17-perm}
  \end{equation}
  where $\mathrm{Perm}(2^n)$ denotes the computational-basis permutation
  matrices.  Their seven-qubit permutation $U_3$ is Clifford-conjugate to the
  Gottesman--Mochon gate (Proposition~5.8), and their Lemma~A.3 shows that
  $U\otimes I$ is not semi-Clifford whenever $U$ is not, so
  $\mathcal{C}_3(n)\setminus\mathrm{SC}(n)\neq\varnothing$ for every
  $n\geq n_{\min}$.  Eq.~\eqref{eq:5a17-perm} does not cover general gates: by
  their Corollary~3.9, every third-level gate has the form $C_L\pi DC_R$ with
  Clifford gates $C_L,C_R$, a permutation gate $\pi\in\mathcal{C}_3(n)$, and a
  diagonal $D$, so a counterexample on five or six qubits would have to arise
  from such a product in which $\pi$ alone is semi-Clifford.  The TQC 2026
  version states in Section~2.1 that the five- and six-qubit cases remain open
  \sourcecite{ref:5a17-hrt26}{HRT26}.

  \item De Silva and Lautsch constructed, in a preliminary preprint of
  10 September 2026, a five-qubit gate $V$, a product of two singly controlled
  Clifford gates, with (Theorem~8.1)
  \begin{equation}
    V\in\mathcal{C}_5(5)\setminus\mathcal{C}_4(5),
    \qquad
    V\ \text{is not generalized semi-Clifford}.
    \label{eq:5a17-fifth}
  \end{equation}
  Tensoring with identities gives such gates in $\mathcal{C}_k(n)$ for all
  $n\geq5$ and $k\geq5$, and $V^\dagger$ lies in no level of the hierarchy
  (Theorem~8.2) \sourcecite{ref:5a17-dsl26}{dSL26}.  Since every third-level
  gate is generalized semi-Clifford \sourcecite{ref:5a17-bs10}{BS10} and
  $V\notin\mathcal{C}_4(5)\supseteq\mathcal{C}_3(5)$,
  Eq.~\eqref{eq:5a17-fifth} has no bearing on $n_{\min}$.
\end{itemize}

\subsection*{References}

\begin{enumerate}[label={},leftmargin=4.8em,labelsep=0.5em]
  \footnotesize
  \item[\textup{[ZCC08]}]\label{ref:5a17-zcc08}
  B. Zeng, X. Chen, and I. L. Chuang, ``Semi-Clifford operations, structure of
  $\mathcal{C}_k$ hierarchy, and gate complexity for fault-tolerant quantum
  computation,'' \emph{Physical Review A} \textbf{77}, 042313 (2008).
  \href{https://doi.org/10.1103/PhysRevA.77.042313}{doi:10.1103/PhysRevA.77.042313};
  \href{https://arxiv.org/abs/0712.2084}{arXiv:0712.2084}.

  \item[\textup{[BS10]}]\label{ref:5a17-bs10}
  S. Beigi and P. W. Shor, ``$\mathcal{C}_3$, semi-Clifford and generalized
  semi-Clifford operations,'' \emph{Quantum Information and Computation}
  \textbf{10}, 41--59 (2010).
  \href{https://doi.org/10.26421/QIC10.1-2-4}{doi:10.26421/QIC10.1-2-4};
  \href{https://arxiv.org/abs/0810.5108}{arXiv:0810.5108}.

  \item[\textup{[AC25]}]\label{ref:5a17-ac25}
  J. T. Anderson and A. Connelly, ``Affine equivalence in the Clifford
  hierarchy,'' arXiv preprint (2025).
  \href{https://doi.org/10.48550/arXiv.2507.14370}{doi:10.48550/arXiv.2507.14370};
  \href{https://arxiv.org/abs/2507.14370}{arXiv:2507.14370}.

  \item[\textup{[HRT26]}]\label{ref:5a17-hrt26}
  Z. He, L. Robitaille, and X. Tan, ``Characterization of permutation gates in
  the third level of the Clifford hierarchy,'' in \emph{21st Conference on the
  Theory of Quantum Computation, Communication and Cryptography (TQC 2026)},
  LIPIcs \textbf{389}, 2:1--2:17 (2026); full version arXiv:2510.04993.
  \href{https://doi.org/10.4230/LIPIcs.TQC.2026.2}{doi:10.4230/LIPIcs.TQC.2026.2};
  \href{https://arxiv.org/abs/2510.04993}{arXiv:2510.04993}.

  \item[\textup{[dSL26]}]\label{ref:5a17-dsl26}
  N. de Silva and O. Lautsch, ``The generalised semi-Clifford conjecture is
  false,'' arXiv preprint, preliminary draft (2026).
  \href{https://doi.org/10.48550/arXiv.2609.11903}{doi:10.48550/arXiv.2609.11903};
  \href{https://arxiv.org/abs/2609.11903}{arXiv:2609.11903}.
\end{enumerate}

\subsection*{Comment}

The unresolved cases are $n=5$ and $n=6$: it is unknown whether
$\mathcal{C}_3(5)$ or $\mathcal{C}_3(6)$ contains a gate that is not
semi-Clifford.  A proof that all gates in both are semi-Clifford would give
$n_{\min}=7$, while one counterexample on five or six qubits would fix
$n_{\min}$ at the smallest such size.  The permutation-gate result in
Eq.~\eqref{eq:5a17-perm} does not exclude products of third-level permutations
with non-Clifford diagonal gates, and the fifth-level gate in
Eq.~\eqref{eq:5a17-fifth} does not supply a third-level counterexample.  This
is a finite-size question; the unrestricted conjecture that all third-level
gates are semi-Clifford is already refuted by Eq.~\eqref{eq:5a17-seven}.  The
four-qubit input in Eq.~\eqref{eq:5a17-four} is a computer-assisted result in a
preprint.  Literature checked through 15 September 2026; no resolution of the
five- or six-qubit case was found.  The zoo's record
``\href{https://qiqc-op.com/problem/op_cdd1f718ccfbccf6/}{Semi-Clifford structure of two-qudit hierarchy gates above the third level}''
asks the analogous semi-Clifford question for two qudits of odd prime
dimension at levels four and above.

\subsection*{Field}

Quantum algorithm

\subsection*{Topic}

Clifford hierarchy

\subsection*{ID}

\texttt{op\_5a17b19b8adeb191}
